% SecondWave --- rapport de feedback développeur
% Compilation :  tectonic FEEDBACK.tex
\documentclass[10pt,a4paper]{article}

\usepackage[T1]{fontenc}
\usepackage{lmodern}

\usepackage[french]{babel}
\usepackage[a4paper,top=1.25cm,bottom=1.1cm,left=1.6cm,right=1.6cm,footskip=0.6cm]{geometry}
\usepackage{booktabs}
\usepackage{tabularx}
\usepackage{xcolor}
\usepackage{enumitem}
\usepackage{titlesec}
\usepackage{fancyvrb}
\usepackage{microtype}
\usepackage[hidelinks]{hyperref}

\definecolor{accent}{HTML}{0D5C5C}
\definecolor{rule}{HTML}{C3CDD3}
\definecolor{shade}{HTML}{F0F3F4}
\definecolor{ink2}{HTML}{3A4854}

% --- titres serrés ---
\titlespacing*{\section}{0pt}{6pt plus 1pt minus 1pt}{2pt}
\titlespacing*{\subsection}{0pt}{4pt plus 1pt minus 1pt}{1pt}
\titleformat{\section}{\large\bfseries\color{accent}}{\thesection\ \textperiodcentered\ }{0pt}{}
\titleformat{\subsection}{\normalsize\bfseries\color{ink2}}{}{0pt}{}

\setlist[itemize]{leftmargin=1.1em,itemsep=1pt,topsep=2pt,parsep=0pt}
\setlist[enumerate]{leftmargin=1.4em,itemsep=1pt,topsep=2pt,parsep=0pt}
\setlength{\parskip}{2pt}
\linespread{0.97}\selectfont
\setlength{\emergencystretch}{2em}
\setlength{\parindent}{0pt}
\renewcommand{\arraystretch}{1.02}

\newcommand{\cat}[2]{{\small\textsf{catégorie} \texttt{#1} \textperiodcentered\ \textsf{lib} #2}\par\vspace{2pt}}
\newcommand{\tc}[1]{\texttt{#1}}

\DefineVerbatimEnvironment{code}{Verbatim}%
  {fontsize=\footnotesize,xleftmargin=6pt,frame=leftline,framerule=1pt,
   rulecolor=\color{accent},framesep=6pt,baselinestretch=0.98}

\sloppy
\begin{document}

% ------------------------------------------------------------------------------------------------------------------------------------------------
\begin{center}
{\LARGE\bfseries Rapport de feedback développeur}\\[2pt]
{\large SecondWave \textperiodcentered\ équipe BFT-PARIS-26}
\end{center}
\vspace{2pt}
{\color{rule}\hrule height 0.6pt}
\vspace{5pt}

{\small
\textbf{Track} 2 --- Lending Protocol \quad \textbf{Flavour} Loaded (Permissioned Domains + Credentials)\\
\textbf{Environnement} Devnet public, \tc{wss://s.devnet.rippletest.net:51233} \textperiodcentered\ rippled 3.4.0-rc5
(\tc{LendingProtocolV1\_1}, \tc{BatchV1\_1}, \tc{fixCleanup3\_4\_0} actifs)\\
\textbf{Libs} \tc{xrpl.js@5.2.0-beta.1} \textperiodcentered\ \tc{ripple-binary-codec@2.11.0} \textperiodcentered\ Node $\geq$ 20
\par}
\vspace{4pt}

\textbf{Le projet} --- un marché secondaire de parts de vault verrouillées. Pendant la phase
Investment d'un vault \emph{closed-ended}, \tc{VaultWithdraw} est bloqué : un déposant qui doit
sortir n'a aucune issue. Les parts étant un MPT, on les lui fait vendre à un autre membre du
fonds, avec un règlement atomique et une analyse du risque à l'achat. Environ 420 cas exécutés
sur la chaîne ; chaque point est rejouable par un script du dépôt. Versions longues :
\tc{probes/FRICTIONS.md} et \tc{probes-marche/FRICTIONS.md}.

% ------------------------------------------------------------------------------------------------------------------------------------------------
\section{Un \texttt{Batch} ne dit pas s'il a produit un effet}
\cat{UX}{rippled 3.4.0-rc5}

\subsection{Le problème}
Précisons d'abord ce qui \textbf{n'est pas} le problème : \tc{tesSUCCESS} sur un
\tc{tfAllOrNothing} dont rien n'a appliqué est défendable --- le drapeau promet \og tout ou
rien\fg{}, et le protocole a fait respecter l'invariant.

Le problème est que \textbf{les deux issues sont indiscernables}. Réponse complète du serveur
sur un \tc{Batch} qui n'a rien fait :

\begin{code}
accepted: true | applied: true | engine_result: tesSUCCESS
engine_result_message: "The transaction was applied."
meta.AffectedNodes: [ 1 noeud ]   <- Balance -150 drops, Sequence +1
\end{code}

Quatre champs qui disent \og ça a marché\fg{}, dont un \tc{applied: true} alors que rien n'a été
appliqué. Le seul n\oe{}ud de métadonnée est la ponction de frais. Et \tc{tx\_json} restitue les
\tc{RawTransactions} --- c'est-à-dire \textbf{ce qu'on a demandé, pas ce qui s'est passé}.

Sur huit \tc{Batch} rejoués depuis le ledger, cinq ayant déplacé des parts et trois n'ayant
rien fait, \textbf{une seule combinaison apparaît} : \tc{tesSUCCESS} / 1 n\oe{}ud. XLS-56 spécifie
pourtant un \tc{BatchExecutions} dans la métadonnée, exactement pour distinguer les deux cas.
Il est absent.

\subsection{Ce qu'on a vécu}
On a passé une demi-journée à croire notre code cassé. L'échange renvoyait \tc{tesSUCCESS},
puis les soldes ne bougeaient pas --- mais pas toujours. Ce n'est pas aléatoire : chaque jambe
est évaluée sur ses propres préconditions, et une seule qui échoue annule tout. Quatre causes
distinctes chez nous, toutes invisibles depuis la réponse :

\vspace{2pt}
{\small\begin{tabularx}{\linewidth}{@{}Xl@{}}
\toprule
\textbf{Cause réelle de l'échec} & \textbf{Code que la jambe aurait rendu seule}\\
\midrule
l'acheteur ne peut pas financer le prix & \tc{tecUNFUNDED\_PAYMENT}\\
l'acheteur n'est pas membre du domaine & \tc{tecNO\_AUTH}\\
l'acheteur n'a pas fait \tc{MPTokenAuthorize} & \tc{tecNO\_AUTH}\\
séquence interne décalée & \tc{terPRE\_SEQ}\\
\bottomrule
\end{tabularx}\par}
\vspace{3pt}

Aggravant : \textbf{\tc{simulate} répond \tc{notImpl} sur un \tc{Batch}} --- ni pré-jouer, ni
constater. Il a fallu deux couches maison : retrouver les jambes par \tc{account\_tx} sur la
plage \tc{[ledgerIndex, ledgerIndex]}, puis comparer les soldes. Environ 150 lignes pour
répondre à \og est-ce que mon échange a eu lieu~?\fg{}.

\subsection{Ce qu'on propose}
\begin{enumerate}
\item \textbf{Livrer le \tc{BatchExecutions} que la spec annonce} --- au minimum, par jambe :
      son index, son hash et son code de résultat.
\item À défaut, \textbf{le strict minimum utile} : l'index de la première jambe en échec et son
      code. Un entier et un code changeraient tout.
\item \textbf{Ne pas écrire \tc{applied: true}} quand aucune jambe n'a appliqué.
\item \textbf{Faire fonctionner \tc{simulate} sur un \tc{Batch}} : c'est la transaction la plus
      difficile à construire et la seule qu'on ne peut pas pré-vérifier.
\end{enumerate}

Le moment est favorable : l'amendement est encore en \tc{BatchV1\_1} sur une \emph{release
candidate}. Changer la métadonnée après activation coûtera bien plus cher.

% ------------------------------------------------------------------------------------------------------------------------------------------------
\section{Trois drapeaux de \texttt{Batch} sur quatre livrent l'actif sans encaisser le prix}
\cat{UX}{rippled 3.4.0-rc5}

\subsection{Le problème}
Le même échange --- parts contre XRP, jambe du prix volontairement impayable --- rejoué sous
chaque drapeau :

\vspace{2pt}
{\small\begin{tabularx}{\linewidth}{@{}Xlll@{}}
\toprule
\textbf{Drapeau} & \textbf{Annoncé} & \textbf{Parts déplacées} & \textbf{Prix payé}\\
\midrule
\tc{tfAllOrNothing} & \tc{tesSUCCESS} & \textbf{0} & 0\\
\tc{tfOnlyOne} \textperiodcentered\ \tc{tfUntilFailure} \textperiodcentered\ \tc{tfIndependent}
  & \tc{tesSUCCESS} & \textbf{300 000} & \textbf{0}\\
\bottomrule
\end{tabularx}\par}
\vspace{3pt}

C'est cohérent avec la sémantique de chaque mode. Mais appliqué à un échange de valeur, trois
modes sur quatre font \textbf{livrer sans encaisser} --- et les quatre annoncent le succès.

\subsection{Ce qu'on a vécu}
On a découvert ça en tabulant les quatre drapeaux par curiosité, pas en cherchant un bug.
L'oubli est impossible --- \tc{Flags: 0} est refusé par \tc{temINVALID\_FLAG}, et le SDK refuse
même de co-signer un \tc{Batch} sans le champ. Le risque est donc \textbf{le mauvais choix},
fait en connaissance de cause. Combiné au point 1, il est \textbf{indétectable} : même code de
retour, même métadonnée, et l'actif est parti.

\subsection{Ce qu'on propose}
Une mise en garde explicite dans XLS-56 pour les échanges de valeur, et un helper de swap dans
les SDK qui impose \tc{tfAllOrNothing}. Le drapeau qui protège devrait être celui qu'on obtient
sans y penser.

% ------------------------------------------------------------------------------------------------------------------------------------------------
\section{Aucune primitive de lecture côté XLS-66}
\cat{missing primitive}{rippled 3.4.0-rc5}

\subsection{Le problème}
\tc{vault\_info} existe. \textbf{Ni \tc{loan\_info}, ni \tc{loan\_broker\_info}, ni la liste des
détenteurs de parts.} Reconstituer l'état d'un vault impose une traversée à trois niveaux,
chacun sur un pseudo-compte différent :

\begin{code}
vault_info(vaultId)
  -> account_objects(pseudo-compte du vault,      type: 'loan_broker')
    -> account_objects(pseudo-compte du courtier, type: 'loan')
\end{code}

Et pour savoir qui détient les parts, aucune commande : il faut \textbf{rejouer tout
l'historique} du pseudo-compte via \tc{account\_tx} et reconstruire les soldes depuis la
métadonnée de chaque transaction.

\subsection{Ce qu'on a vécu}
Notre produit est une aide à la décision : il doit afficher l'état d'un vault avant qu'on
achète une position. \textbf{Six appels RPC} pour un vault à 3 courtiers et 3 prêts, plus un
parcours paginé complet pour les détenteurs. C'est la partie du code qui nous a demandé le plus
de travail, et elle ne produit aucune valeur métier --- elle reconstitue ce que le serveur sait
déjà.

Un point positif : le rejeu est \textbf{auto-vérifiant} --- la somme reconstruite retombe
exactement sur \tc{shares.OutstandingAmount}.

\subsection{Ce qu'on propose}
\tc{loan\_broker\_info} et \tc{loan\_info} sur le modèle de \tc{vault\_info}, et surtout une
commande qui liste les détenteurs d'une émission MPT. Le rejeu d'historique devrait être un cas
limite, pas le chemin normal.

% ------------------------------------------------------------------------------------------------------------------------------------------------
\section{\texttt{LendingProtocolV1\_1} est actif et n'est documenté nulle part}
\cat{documentation/tutorials}{rippled 3.4.0-rc5}

\subsection{Le problème}
\tc{VaultKind}, \tc{SubscriptionDate}, \tc{RedemptionDate} et \tc{LEVersion} existent sur le
ledger et sont connus du codec. \textbf{Aucune spécification publiée ne les décrit.} Le modèle
en trois phases, l'obligation d'un vault closed-ended pour créer un courtier, les bornes des
dates : tout découvert par essai-erreur.

Aggravant : \textbf{côté serveur, toute violation de \tc{VaultCreate} sort en \tc{temMALFORMED}
nu}, sans jamais nommer le champ. Les messages utiles --- \og RedemptionDate $-$ SubscriptionDate
must be within [180, 946708560) seconds\fg{} --- n'existent \textbf{que dans la validation client de
\tc{xrpl.js}}. Un intégrateur en Python, en Java, ou qui signe à la main, n'a rien.

\subsection{Ce qu'on a vécu}
Le cas qui nous a le plus coûté : tout \tc{LoanSet} avec \tc{GracePeriod < 60} ou
\tc{> PaymentInterval} échoue en \tc{temINVALID} \og The transaction is ill-formed\fg{}, sans
indication de champ. \textbf{Une session de debug entière}, trois hypothèses fausses, puis un
balayage systématique pour isoler la règle --- qui n'est écrite nulle part.

\vspace{2pt}
{\small\begin{tabularx}{\linewidth}{@{}llX@{}}
\toprule
\textbf{Envoyé} & \textbf{Reçu} & \textbf{Signifiait réellement}\\
\midrule
\tc{GracePeriod: 59} & \tc{temINVALID} & plancher non documenté à 60 s\\
\tc{LoanBrokerSet} sur vault open-ended & \tc{tecNO\_PERMISSION} & V1\_1 exige un closed-ended\\
échéance proche de \tc{RedemptionDate} & \tc{tecNO\_PERMISSION} & marge réelle $\sim$90 s en temps ledger\\
\tc{StartDate} dans \tc{LoanSet} & \emph{disallowed location} & champ dérivé sur l'objet \tc{Loan}\\
\tc{Data} dans \tc{LoanSet} & \tc{tesSUCCESS} & \textbf{accepté puis silencieusement jeté}\\
\bottomrule
\end{tabularx}\par}
\vspace{3pt}

\subsection{Ce qu'on propose}
Publier la spec de V1\_1, même à l'état de brouillon --- un amendement actif sur un réseau public
devrait avoir une page. Et \textbf{remonter côté serveur les messages qui n'existent aujourd'hui
que dans le SDK JavaScript} : c'est là qu'ils servent à tout le monde, quel que soit le langage.

% ------------------------------------------------------------------------------------------------------------------------------------------------
\section{La co-signature \texttt{LoanSet} était invalide --- corrigé en \texttt{5.2.0-beta.1}}
\cat{client libraries}{\tc{xrpl.js} 5.2.0-beta.0 $\rightarrow$ beta.1}

\subsection{Le problème}
\tc{signLoanSetByCounterparty} signait avec le préfixe historique \tc{STX\textbackslash0} alors
que \tc{fixCleanup3\_4\_0}, actif sur ce Devnet, exige \tc{CPT\textbackslash0} en simple
signature et \tc{CPM\textbackslash0} en multisig. Le n\oe{}ud refusait aux \emph{local checks} :
\tc{fails local checks: Counterparty: Invalid signature.}

\textbf{Aucun \tc{LoanSet} n'était soumettable depuis le SDK publié} --- la transaction centrale
de XLS-66, et l'exemple officiel \tc{createLoan.js} avec elle. Les bons encodeurs existaient
déjà dans \tc{ripple-binary-codec@2.11.0}, avec l'amendement nommé en commentaire ; ils
n'étaient jamais appelés. Et \tc{.ci-config/xrpld.cfg} s'arrêtant à \tc{fixCleanup3\_2\_0}, la
CI ne pouvait pas voir la casse.

\subsection{Ce qu'on a vécu}
Deux heures à chercher une erreur de notre côté, jusqu'au moment où le même payload signé avec
\tc{encodeForSigningCounterparty} est passé. Sans prêt, un vault n'a ni rendement ni risque ---
\textbf{tout notre projet était bloqué} derrière cette signature. \tc{5.2.0-beta.1} route
désormais la signature par une table d'encodeurs indexée sur le rôle ; revérifié sur la chaîne :
\tc{tesSUCCESS}, objet \tc{Loan} créé.

\subsection{Ce qu'on propose}
Deux points restent ouverts. \textbf{Ajouter \tc{fixCleanup3\_4\_0} à la configuration CI} de
\tc{xrpl.js} --- sans ça la régression peut revenir sans qu'aucun test ne bronche, et c'est
exactement ce qui a laissé passer le bug. Et \textbf{\tc{xrpl-py} (\tc{counterparty\_signer.py})
porte toujours le même défaut}, son codec ayant d'abord besoin des deux préfixes.

% ------------------------------------------------------------------------------------------------------------------------------------------------
\section{Les MPT dans \texttt{Escrow} fonctionnent et ne sont documentés nulle part}
\cat{documentation/tutorials}{rippled 3.4.0-rc5}

\subsection{Le problème}
En cherchant un rail de règlement alternatif au \tc{Batch}, on a testé les quatre candidats.
Trois sont fermés et le disent clairement --- \tc{CheckCreate} $\rightarrow$ \tc{invalid SendMax},
\tc{PaymentChannelCreate} $\rightarrow$ \tc{Amount must be a string}, \tc{OfferCreate}
$\rightarrow$ \tc{temDISABLED}. Le quatrième marche de bout en bout, \textbf{et rien ne le dit} :

\begin{itemize}
\item \tc{EscrowCreate} avec des parts de vault $\rightarrow$ \tc{tesSUCCESS}, parts débitées à la création ;
\item \tc{EscrowFinish} \textbf{livre réellement}, et \textbf{crée l'objet \tc{MPToken} du
      destinataire tout seul} --- là où un \tc{Payment} exige un \tc{MPTokenAuthorize} préalable ;
\item \tc{EscrowCancel} restitue l'intégralité, et le contrôle de domaine s'applique
      \textbf{à la création comme au dénouement} ;
\item deux escrows partageant une même \tc{Condition} forment un échange atomique \textbf{sans
      \tc{Batch}} --- mesuré, 4 $\times$ \tc{tesSUCCESS}.
\end{itemize}

\subsection{Ce qu'on a vécu}
Trouvé par hasard, en testant les rails un par un pour documenter ce qui est fermé. Capacité
utile et invisible : la documentation d'\tc{Escrow} ne mentionne pas les MPT, et rien dans
XLS-33 ne renvoie vers \tc{Escrow}.

\textbf{Un cas limite à signaler} : un escrow \textbf{sans \tc{CancelAfter}} dont le destinataire
perd son credential devient irrécupérable --- \tc{EscrowFinish} rend \tc{tecNO\_AUTH},
\tc{EscrowCancel} rend \tc{tecNO\_PERMISSION}, les deux sorties sont fermées. Reproduit de bout
en bout ; nous le signalons sans savoir si c'est voulu.

\subsection{Ce qu'on propose}
Documenter la compatibilité MPT $\times$ \tc{Escrow}, et l'asymétrie d'autorisation : recevoir un
MPT par \tc{Payment} exige \tc{MPTokenAuthorize}, par \tc{EscrowFinish} non. C'est pour
contourner la première que notre \tc{Batch} compte \textbf{trois jambes au lieu de deux}.

% ------------------------------------------------------------------------------------------------------------------------------------------------
\vspace{4pt}
{\color{rule}\hrule height 0.6pt}
\vspace{5pt}

\textbf{\color{accent}Ce qui nous a rassurés}\par
\textbf{Personne ne peut enfermer un déposant} : les dates ne sont pas des champs de
\tc{VaultSet}, et un membre exclu d'un domaine peut toujours retirer. \textbf{Le pseudo-compte
est étanche} --- \tc{Payment}, \tc{EscrowCreate}, \tc{CheckCreate}, \tc{PaymentChannelCreate}
$\rightarrow$ \tc{tecNO\_PERMISSION} ; vault dans un \tc{Batch} $\rightarrow$
\tc{temINVALID\_INNER\_BATCH}. \textbf{Aucune fuite d'arrondi} sur 145 aller-retours en XRP et
en IOU. \textbf{Les \tc{BatchSigners} couvrent les jambes internes} : un vendeur qui réécrit le
prix après signature de l'acheteur est rejeté, et sa séquence interdit la double-vente. Enfin
\textbf{\tc{tfAllOrNothing} tient} dans tous les cas d'échec testés --- ce qui manque n'est pas
la garantie, c'est le moyen de la constater.

\vspace{3pt}
\textbf{\color{accent}Une question ouverte}\par
\tc{OfferCreate} avec un MPT répond \tc{temDISABLED} : le code existe, l'amendement dort.
\textbf{Le carnet d'ordres natif connaîtra-t-il les Permissioned Domains ?} Nos parts sont
gatées ; un carnet sans permission apparierait des ordres dont le règlement serait refusé. Selon
la réponse, notre marché garde tout son sens --- ou devient inutile.

\end{document}
